\section{Decoder 预训练：GPT 系列与大语言模型}

\subsection{Decoder 的预训练与微调}

Decoder-only 模型的预训练目标就是标准的\textbf{语言建模}：

$$\mathcal{L}(\theta) = -\sum_{t=1}^{T} \log p_\theta(w_t | w_1, \ldots, w_{t-1})$$

\begin{figure}[H]
\centering
\includegraphics[width=0.95\textwidth]{slides-images/slide-045.jpg}
\caption{Decoder 预训练与微调：预训练时用语言建模目标；微调时取最后一个 token 的隐状态接线性分类器，梯度反传整个网络\protect\footnotemark}
\end{figure}
\footnotetext{Slides 第46页。}

微调方式与 BERT 类似：取序列最后一个 token 的隐状态 $h_T$，接一个线性层 $y \sim Ah_T + b$ 进行分类。

\subsection{GPT-1：开创性的预训练 Decoder}

\textbf{GPT}（Generative Pre-Training，2018，OpenAI）是第一个大规模预训练的 Decoder-only 模型：

\begin{itemize}
    \item 117M 参数，12 层，768 维隐层
    \item 约 40,000 词的 BPE 词表
    \item 在 BookCorpus 上训练
\end{itemize}

\begin{figure}[H]
\centering
\includegraphics[width=0.95\textwidth]{slides-images/slide-047.jpg}
\caption{GPT 用于自然语言推断（NLI）：将前提和假设用分隔符连接，作为一个序列输入 Decoder，取最后一个 token 的表示进行分类\protect\footnotemark}
\end{figure}
\footnotetext{Slides 第48页。}

GPT 在多项 NLP 基准上取得了强劲表现。BERT 紧随其后发布，凭借双向上下文在分类任务上略胜一筹，但 GPT 的生成能力是 BERT 无法匹敌的。

\subsection{GPT-2：惊人的生成能力}

GPT-2（2019）将模型规模提升至 \textbf{1.5B 参数}，训练数据增至约 90 亿词。

\begin{figure}[H]
\centering
\includegraphics[width=0.95\textwidth]{slides-images/slide-050.jpg}
\caption{GPT-2 的生成示例：给定一段关于``科学家发现独角兽''的 prompt（人类编写），GPT-2 生成了一段高度连贯且富有创意的续写\protect\footnotemark}
\end{figure}
\footnotetext{Slides 第51页。来自 Radford et al., 2019。}

GPT-2 的生成质量在当时令人震惊，它能够产生\textbf{多段落、主题连贯}的文本，且 1.5B 的规模仍然可以在小型 GPU 上进行微调。

\subsection{GPT-3：涌现的上下文学习能力}

GPT-3（2020）实现了参数规模的巨大飞跃：

\begin{table}[H]
\centering
\begin{tabular}{lccc}
\toprule
\textbf{模型} & \textbf{参数量} & \textbf{训练数据} & \textbf{年份} \\
\midrule
GPT-1 & 117M & BookCorpus & 2018 \\
GPT-2 & 1.5B & WebText（9B 词） & 2019 \\
GPT-3 & 175B & 300B 词 & 2020 \\
\bottomrule
\end{tabular}
\caption{GPT 系列模型的规模演进}
\end{table}

GPT-3 最令人惊讶的不是它的微调性能，而是它展现出了一种全新的能力——\textbf{上下文学习}（In-Context Learning, ICL）。

\begin{figure}[H]
\centering
\includegraphics[width=0.95\textwidth]{slides-images/slide-052.jpg}
\caption{上下文学习示例：在 prompt 中给出几个英法翻译的例子（thanks $\rightarrow$ merci, hello $\rightarrow$ bonjour），模型能够继续这个``模式''，为新词生成翻译\protect\footnotemark}
\end{figure}
\footnotetext{Slides 第53页。来自 Brown et al., 2020。}

\begin{importantbox}{上下文学习（In-Context Learning）}
无需任何梯度更新或微调，仅通过在 prompt 中给出几个\textbf{示例}（demonstrations），模型就能``学会''执行新任务。这完全是通过\textbf{模式匹配}实现的：模型识别出 prompt 中的输入-输出模式，并将其应用于新的输入。

GPT-3 的上下文学习能力涵盖：
\begin{itemize}
    \item 翻译（给出几对翻译示例）
    \item 算术（给出几个加法示例）
    \item 纠错（给出几个错别字修正示例）
    \item 以及更多...
\end{itemize}
\end{importantbox}

\begin{figure}[H]
\centering
\includegraphics[width=0.95\textwidth]{slides-images/slide-053.jpg}
\caption{更多上下文学习示例：简单算术、纠错、翻译等——这些能力并非通过微调获得，而是在预训练中\textbf{涌现}出来的\protect\footnotemark}
\end{figure}
\footnotetext{Slides 第54页。}

Chris Manning 强调这种能力是``qualitatively new behavior''——从小模型中完全无法预见。GPT-3 只是被训练来预测下一个词，但它\textbf{涌现}出了执行复杂任务的能力。这既令人兴奋，也令人不安，因为我们无法准确预测什么能力会涌现、什么时候会涌现。

\subsection{Scaling Laws 与计算效率}

\begin{figure}[H]
\centering
\includegraphics[width=0.95\textwidth]{slides-images/slide-055.jpg}
\caption{Scaling Laws：测试损失与计算量、数据量、参数量之间存在幂律关系——增大任何一个维度都能可预测地降低损失\protect\footnotemark}
\end{figure}
\footnotetext{Slides 第56页。来自 Kaplan et al., 2020。}

\begin{knowledgebox}{Chinchilla Scaling Laws}
DeepMind 的 Chinchilla 研究（Hoffmann et al., 2022）发现 GPT-3 的参数量与训练数据量之间\textbf{严重不平衡}——GPT-3 ``comically oversized''。Chinchilla 只有 70B 参数（不到 GPT-3 的一半），但在 1.4T token 上训练（GPT-3 只用了 300B），结果性能更优。

核心结论：给定固定计算预算，参数量 $N$ 和训练 token 数 $D$ 应该\textbf{等比例}缩放：$N \propto D$。大约每个参数对应 20 个训练 token。
\end{knowledgebox}

\subsection{Chain-of-Thought 提示}

\begin{figure}[H]
\centering
\includegraphics[width=0.95\textwidth]{slides-images/slide-058.jpg}
\caption{Chain-of-Thought Prompting：左边的标准提示直接给出答案（错误）；右边在示例中展示\textbf{推理步骤}，模型模仿这一模式进行逐步推理，得出正确答案\protect\footnotemark}
\end{figure}
\footnotetext{Slides 第59页。来自 Wei et al., 2022。}

\begin{importantbox}{Chain-of-Thought Prompting 的原理}
标准 prompting：问题 $\rightarrow$ 答案

Chain-of-Thought：问题 $\rightarrow$ \textbf{推理步骤} $\rightarrow$ 答案

关键洞察：模型在生成每个 token 时能\textbf{条件于之前生成的所有 token}。通过先生成中间推理步骤，模型：
\begin{itemize}
    \item 获得了额外的``思考空间''（类似草稿纸）
    \item 将复杂问题分解为更简单的子问题
    \item 在每个子步骤中能利用之前步骤的结论作为条件
\end{itemize}
这显著提升了数学推理、逻辑推理等需要多步思考的任务的准确率。
\end{importantbox}

\subsection{为什么 Decoder-only 模型占据主导地位？}

尽管 Encoder 和 Encoder-Decoder 各有优势，当今最大、最强的预训练模型几乎都是 Decoder-only 架构。Chris Manning 指出原因可能包括：

\begin{enumerate}
    \item \textbf{简洁性}：Decoder-only 架构更简单，所有参数共享在一个网络中
    \item \textbf{参数效率}：Encoder-Decoder 需要将参数分配给两个子网络，而 Decoder-only 的所有参数都用于同一个任务
    \item \textbf{统一性}：所有任务（理解+生成）都可以建模为``给定 prompt，生成 completion''
    \item \textbf{涌现能力}：上下文学习等涌现能力在 Decoder-only 模型中表现最为突出
\end{enumerate}

\subsection{本章小结}

Decoder-only 预训练以 GPT 系列为代表，通过简单的语言建模目标学习强大的语言能力。从 GPT-1（117M）到 GPT-3（175B），模型规模的增长带来了质变式的涌现能力，尤其是上下文学习。Chinchilla 揭示了``更大不一定更好''的计算效率规律。Chain-of-Thought prompting 则展示了引导模型进行多步推理的有效策略。Decoder-only 架构已成为当今大语言模型的主流选择。

%% ========== Part 6: 预训练数据 ==========

